\subsection{EQUIVALENCE RELATIONS ON A REGULAR SPACE}

PROPOSITION 14. Let X be a regular space, R a closed equivalence relation on X. Then the graph C of R in X × X is closed.

Let \((a, b)\) be a point of \(X \times X\) in the closure of C, and let V (resp. W) be a closed neighbourhood of \(a\) (resp. a neighbourhood of \(b\) ) in X; then there is a point \((x, y) \in C \cap (V \times W)\) . Since \(x \in V\) , \(y\) belongs to the saturation S of V with respect to R; hence each neighbourhood W of \(b\) meets S. By hypothesis S is closed, and therefore \(b \in S\) . Now let B be the saturation of \(\{b\}\) with respect to R, then each closed neighbourhood V of \(a\) meets B; since by hypothesis B is closed and X is regular, it follows that \(a \in B\) and therefore that \((a, b) \in C\) . This completes the proof.

COROLLARY. On a regular space, every equivalence relation which is both open and closed is Hausdorff.

This follows from Proposition 14 and Proposition 8 of no. 3.

PROPOSITION 15. Let X be a regular space, F a closed subset of X, R the equivalence relation on X obtained by identifying all the points of F {[}in other words, the equivalence relation whose equivalence classes are F (if \(\mathbf{F} \neq \emptyset\) ) and the sets \(\{x\}\) where \(x \in [\mathbb{C}F]\) . Then the quotient space X/R is Hausdorff.

Let M and N be two distinct equivalence classes in X. If each of them consists of a single point in the complement of F, then there exist two disjoint open neighbourhoods of M and N in the Hausdorff subspace \(\complement F\); these are neighbourhoods of M and N in X and are saturated with respect to R. If M = F (so that \(F \neq \varnothing\) ) and \(N = \{b\}\) where \(b \notin F\) , then since X is regular there is an open neighbourhood of b and an open neighbourhood of F which do not intersect; these neighbourhoods are saturated with respect to R, and the proposition is proved.

Note that the quotient space X/R is not necessarily regular (Chapter IX, § 4, Exercise 14).

